\documentclass[11pt]{article}
\usepackage[a4paper,margin=25mm]{geometry}
\usepackage[no-math]{fontspec}
\setmainfont{Latin Modern Roman}
\usepackage{amsmath,amssymb,amsthm,xfrac,xspace,hyperref,graphicx,comment}
\excludecomment{DefTralics}
\providecommand{\C}{}
\newcommand{\ie}{i.e.,\xspace}
\newcommand{\eg}{e.g.,\xspace}
\newcommand{\etc}{etc.\xspace}
\newcommand{\cf}{cf.\xspace}
\newcommand{\resp}{resp.\xspace}
\newcommand{\smaller}{\small}
\newcommand{\Subsection}{\subsection}
\newcommand{\Subsubsection}{\subsubsection}
\newcommand{\skpt}{\smallskip}
\emergencystretch=3em
\input{author-preamble.tex}
\begin{document}
\begin{center}{\Large Stable item 1926}\end{center}
\paragraph{Source and attribution.}
Olivier Benoist, Olivier Wittenberg, \emph{On the integral Hodge conjecture for real varieties, II}, Journal de l’École polytechnique — Mathématiques 7 (2020), publisher version, DOI 10.5802/jep.120. \href{https://jep.centre-mersenne.org/articles/10.5802/jep.120/}{Publisher article}. Authors' text: \href{https://creativecommons.org/licenses/by/4.0/}{CC BY 4.0}.
\paragraph{Editorial problem boundary.}
Prove only Proposition 9.17 below, the explicit failure of the EPT property on a K3 surface over the real Puiseux field \(R=\bigcup_n\mathbb R((t^{1/n}))\), with the usual positive-infinitesimal ordering. The selected surface, double-cover equation, strict-transform line bundle, vanishing Borel--Haefliger class, and full obstruction to a real-point-free divisor representative occur in the original proof. Definition 9.15 of EPT is retained. No assertion over arbitrary non-archimedean fields or integral-Hodge-conjecture application is selected; the rest of this article is conservative author context only.
\paragraph{Difficulty (editorial): H3.}
An explicit branched double-cover K3 construction converts a hypothetical real-point-free divisor into an odd-bidegree invariant section. Normalization over a Puiseux valuation ring and specialization of real spectra force an odd-degree real curve on which the defining polynomial must have an impossible constant sign.
\paragraph{Editorial transcription note.}
Original author language, mathematical content and ordering are unchanged. Publisher front matter is replaced by this independent layout wrapper. Original native body and macros remain editable; bibliography is recovered from the matching cached publisher HTML. Adjacent claims are author context, not additional corpus items. Printed mathematical slips are retained without assistant repair.
\clearpage
\paragraph{Standard background (editorial).}
The Borel--Haefliger class of a real line bundle and basic real-spectrum/normalisation facts are standard prerequisites; the class-map setup is discussed in the authors' first article \cite{BW1}, and the exact real-spectrum specialisation fact used in the proof is cited by the authors to \cite[Prop.\,7.1.21]{bcr}. No proof from the first article is asserted or required as a new core step of this selected construction.
\setcounter{section}{9}\setcounter{subsection}{3}\setcounter{subsubsection}{0}\setcounter{thm}{14}
\input{author-definition.tex}
\setcounter{thm}{16}\setcounter{equation}{3}
\input{author-proof.tex}
\begin{thebibliography}{99}
\bibitem[BCR98]{bcr} J. Bochnak, M. Coste \& M.-F. Roy - Real algebraic geometry , Ergeb. Math. Grenzgeb. (3) , vol. 36 , Springer-Verlag, Berlin, 1998
\bibitem[BW18]{BW1} O. Benoist \& O. Wittenberg - “On the integral Hodge conjecture for real varieties, I” , 2018, à paraître dans Invent. math.
\end{thebibliography}
\end{document}
